\begin{proof}[Proof of \cref{obj:prop:published-class-converse-transfer}]
Fix \(\beta\in(0,1]\).

\begin{enumerate}
\item By \cref{obj:thm:uniform-frontier}, there is a constant \(c>0\), depending only on \(\beta\), such that, simultaneously for every \(n\ge2\) and \(\sigma\in[0,1]\),
\[
c\,r_\beta(n,\sigma)\le R_\beta(n,\sigma),
\qquad
c\,r_\beta(n,\sigma)\le\mathcal L_\beta(n,\sigma).
\]
We use this same constant for both transferred lower bounds.
% lean: CausalSmith.Stat.RdTruesideNoiseFrontier.uniform_frontier

\item Fix \(\sigma\in[0,1]\) and \(P\in\mathcal M_\beta(\sigma)\). Define the embedded law and its accompanying density and mean versions by
\[
Q=P,\qquad
(Z,V^0,V^1,E)=(X,Y^0,Y^1,\varepsilon),
\qquad
a(z)=f(P,z),\qquad v_d(z)=\mu_d(P,z)
\quad(d\in\{0,1\},\ z\in\mathbb R).
\]
The density identity is inherited from \cref{obj:synth_1}. For every measurable set \(F_{\mathrm{score}}\subseteq\mathbb R\), the conditional-mean identity becomes
\[
\begin{aligned}
\int_{\{Z\in F_{\mathrm{score}}\}}V^d\,dQ
&=\int_{F_{\mathrm{score}}\cap[-1,1]}
       \mu_d(P,z)f(P,z)\,dz\\
&=\int_{F_{\mathrm{score}}}\mu_d(P,z)f(P,z)\,dz\\
&=\int_{F_{\mathrm{score}}}v_d(z)a(z)\,dz.
\end{aligned}
\]
The first equality is supplied by \cref{obj:synth_1}; the second follows because \(f(P,z)=0\) outside \([-1,1]\). Thus the chosen \(v_d\) remain conditional-mean versions on the real score space.
% lean: CausalSmith.Stat.RdTruesideNoiseFrontier.benchmark_version_on_real

Choose
\[
r=\frac12.
\]
Since \((-r,r)\subseteq[-1,1]\), the full-interval continuity in \cref{obj:synth_1} gives continuity of \(a\) and both \(v_d\) on \((-r,r)\). Moreover, \cref{obj:ass:density-bounds,obj:ass:mean-bounds} give
\[
a(z)\ge\frac14>0\quad(z\in(-r,r)),
\qquad
0\le\frac14\le v_d(0)\le\frac34\le1
\quad(d\in\{0,1\}).
\]
The Gaussian condition follows from \cref{obj:ass:gaussian}:
\[
E\sim N(0,1).
\]
To verify the required joint independence, define the measurable map
\[
\Phi_{\mathrm{obs}}(z,y_0,y_1)
=
\left(
z,\,
\mathbf1\{z\ge0\}y_1+
\bigl(1-\mathbf1\{z\ge0\}\bigr)y_0
\right),
\qquad
(z,y_0,y_1)\in\mathbb R\times\{0,1\}^2.
\]
Then
\[
(Z,B)=\Phi_{\mathrm{obs}}(X,Y^0,Y^1).
\]
By \cref{obj:ass:error-independence}, independence is preserved under this measurable transformation, so
\[
E\perp(Z,B).
\]
All defining conditions of \(\mathcal G(\sigma)\) therefore hold.
% lean: CausalSmith.Stat.RdTruesideNoiseFrontier.benchmarkEmbedding_mem

The embedding leaves the coordinates used by the observation map unchanged. Consequently,
\[
\vartheta(Q)=v_1(0)-v_0(0)=\theta(P),
\qquad
Q\circ(R,A,B)^{-1}=P_{\mathrm{obs}},
\]
and hence, for every integer \(n\ge0\),
\[
\bigl(Q\circ(R,A,B)^{-1}\bigr)^{\otimes n}
\otimes\operatorname{Unif}[0,1]
=
P_{\mathrm{obs}}^{\otimes n}
\otimes\operatorname{Unif}[0,1]
=
\mathbb P_{P,n}.
\]
At \(n=0\), the sample factor is the empty product probability law, so the identity also covers that case.
% lean: CausalSmith.Stat.RdTruesideNoiseFrontier.benchmarkEmbedding_preserves

\item Fix \(n\ge2\) and \(\sigma\in[0,1]\). For every measurable estimate \(T\), target and experiment preservation give
\[
\sup_{P\in\mathcal M_\beta(\sigma)}
\mathbb E_{P,n}|T-\theta(P)|
\le
\sup_{Q\in\mathcal G(\sigma)}
\mathbb E_Q|T-\vartheta(Q)|.
\]
Taking infima over the common estimator class gives the corresponding comparison of extended nonnegative minimax objectives.

To justify the comparison after conversion to real values, note that every \(Q\in\mathcal G(\sigma)\) satisfies
\[
-1\le v_1(0)-v_0(0)=\vartheta(Q)\le1.
\]
Define the constant estimate by
\[
T_{\mathrm{zero}}(S_n,U)=0.
\]
Since the randomized experiments are probability laws,
\[
\sup_{Q\in\mathcal G(\sigma)}
\mathbb E_Q|T_{\mathrm{zero}}-\vartheta(Q)|
\le1.
\]
Thus the larger-class extended minimax risk is finite, and conversion to real values preserves the preceding inequality:
\[
R_\beta(n,\sigma)\le\mathcal R_{\mathcal G}(n,\sigma).
\]
% lean: CausalSmith.Stat.RdTruesideNoiseFrontier.benchmark_risk_le_gaussian

For interval length, define the constant interval procedure by
\[
I_{\mathrm{full}}(S_n,U)=[-1,1].
\]
The same target bounds imply, for every \(Q\in\mathcal G(\sigma)\),
\[
Q\{\vartheta(Q)\in I_{\mathrm{full}}\}=1,
\qquad
\mathbb E_Q\operatorname{length}(I_{\mathrm{full}})=2.
\]
Hence the larger-class honest interval class is nonempty and its extended minimax expected length is at most \(2\).

Now let \(I=[\ell,\upsilon]\) be any interval procedure honest on \(\mathcal G(\sigma)\). For every benchmark law \(P\), its embedded law \(Q\) satisfies
\[
\mathbb P_{P,n}\{\theta(P)\in I\}
=
Q\{\vartheta(Q)\in I\}
\ge0.9.
\]
Thus \(I\) belongs to the benchmark honest interval class of \cref{obj:def:honest-class}. Moreover,
\[
\mathcal L_\beta(n,\sigma)
\le
\sup_{P\in\mathcal M_\beta(\sigma)}
\mathbb E_{P,n}[\upsilon-\ell]
\le
\sup_{Q\in\mathcal G(\sigma)}
\mathbb E_Q[\upsilon-\ell].
\]
The second inequality follows from experiment preservation. Taking the infimum over intervals honest on \(\mathcal G(\sigma)\) proves the extended minimax comparison; the finite bound \(2\) permits its conversion to real values. Therefore
\[
\mathcal L_\beta(n,\sigma)\le\mathcal L_{\mathcal G}(n,\sigma).
\]
Combining these comparisons with the lower bounds already obtained yields
\[
c\,r_\beta(n,\sigma)
\le R_\beta(n,\sigma)
\le\mathcal R_{\mathcal G}(n,\sigma),
\qquad
c\,r_\beta(n,\sigma)
\le\mathcal L_\beta(n,\sigma)
\le\mathcal L_{\mathcal G}(n,\sigma).
\]
% lean: CausalSmith.Stat.RdTruesideNoiseFrontier.benchmark_length_le_gaussian

\item Finally, fix \(\sigma\in(0,1]\) and \(Q\in\mathcal G(\sigma)\). Its joint independence condition, followed by the measurable scaling map \(e\mapsto\sigma e\), gives
\[
E\perp(Z,B),
\qquad
\sigma E\perp(Z,B).
\]
These supply the standardized and scaled error independence in the stated Assumption 1Y condition. The local density condition supplies a radius \(r>0\) for which
\[
a(z)>0\qquad(z\in(-r,r)),
\]
which is Assumption 2C. Finally, scaling the standard Gaussian error gives
\[
N(0,1)\circ(e\mapsto\sigma e)^{-1}=N(0,\sigma^2),
\qquad
\sigma E\sim N(0,\sigma^2).
\]
This is Assumption 7, completing the three asserted conditions.
% lean: CausalSmith.Stat.RdTruesideNoiseFrontier.gaussian_published_hypotheses
\end{enumerate}
\end{proof}
